\section{C-TreePO Objective, Bias, and Sampling}
\label{sec:v8-objective}

\subsection{The C-TreePO Objective}

Section~\ref{sec:v8-ctree-math} gives the structural goal: learn one state map
\(g\) whose node states can replace the raw spans they represent for
oracle-compatible readouts. C-TreePO turns that structural goal into an
optimization problem for the same learned \(g\). The root should be useful for
the document-level task, and the internal states should behave like valid
C-Tree states.

The objective therefore has two channels. The root channel fits the final
document-level target. The local-law channel asks whether the same \(g\)
behaves correctly at leaves, on re-entry, and at merges. Start with the clean
oracle-available version, where \(f^\ast\) can be queried at every relevant
node. This is not the usual training regime, but it is the target objective.

Let \(L_{\mathrm{root}}(g)\) be the full-document root loss, measured against
the available document-level expert target. For the local channel, attach one
node loss to each \(v\in V(T)\). This node loss collects the local checks that
apply at \(v\): C1 for a leaf, C2 for the produced state, and C3 for an
internal merge. If every check could be measured by the target oracle, the
quantity we would calculate at node \(v\) is
\[
  \underbrace{\ell_v^\ast}_{\text{oracle loss we want}}
  =
  \underbrace{\ell_{\mathrm{law}}(v;f^\ast,g)}_{\text{C1/C2/C3 discrepancy through }f^\ast}.
\]
The star marks the target quantity: the local-law loss as it would be measured
by \(f^\ast\), not by the cheaper proxy used during dense training.

Let \(d(v)\) be the depth of node \(v\), with the root at depth \(0\), and let
\(\gamma\ge 0\) be a declared depth-discount weight. If every relevant node
could be evaluated by \(f^\ast\), the local-law part of the objective would be
\begin{equation}
\label{eq:v8-local-law-estimand}
  L_{\mathrm{law}}^\ast(g,T)
  =
  \sum_{v\in V(T)}\gamma^{d(v)}\ell_v^\ast .
\end{equation}
The theorem-facing objective is then
\begin{equation}
\label{eq:v8-main-objective}
  J_\lambda^\ast(g)
  =
  (1-\lambda)L_{\mathrm{root}}(g)
  +
  \lambda L_{\mathrm{law}}^\ast(g,T),
  \qquad 0\le\lambda\le 1 .
\end{equation}
Here \(\lambda\) is the analyst-chosen local-law share. This is the target
objective in principle: root supervision plus the oracle local-law losses that
would make the exact theorem relevant to the fitted \(g\).

\subsection{Why the Objective Becomes an Estimation Problem}

The statistical problem is not that we want an alternate target. It is that
\(L_{\mathrm{law}}^\ast\) is usually not fully observed. At scale, \(f^\ast\)
is sparse, while the local laws create many node-level checks.

We usually observe three different things. First, the root channel has
document-level labels or expert targets for some documents. Second, the
local-law channel has cheap proxy measurements at many nodes. Third, a smaller
sample of nodes is sent to \(f^\ast\), so that the proxy can be checked against
the oracle.

The dense local quantity is the same C1/C2/C3 comparison measured through the
fitted proxy readout \(\widehat f\):
\[
  \underbrace{\widehat\ell_v}_{\text{proxy loss we observe densely}}
  =
  \underbrace{\ell_{\mathrm{law}}(v;\widehat f,g)}_{\text{same local check through }\widehat f}.
\]
The proxy is useful because it gives local training signal everywhere. It is
not an independent target. It supplies the dense local geometry that constrains
\(g\) between sparse oracle calls. The sparse oracle calls then tell us how
the dense proxy surface differs from the target oracle surface, and therefore
how much bias must be corrected in the local-law channel.

The difference between the dense proxy loss and the oracle loss is
\[
  b_v=\widehat\ell_v-\ell_v^\ast .
\]
This \(b_v\) is the proxy bias for node \(v\). If it were always zero, the
proxy channel would already be the oracle local-law channel. In general it is
not zero, so sampled oracle calls are used to estimate and correct it.

Let \(R_v\in\{0,1\}\) indicate whether node \(v\) is sampled and evaluated
with \(f^\ast\). Let \(\pi_v>0\) be the logged marginal probability that the
audit design would sample node \(v\). Thus \(R_v/\pi_v\) is zero for unsampled
nodes and \(1/\pi_v\) for sampled oracle nodes.

\subsection{Correcting the Local-Law Channel}

The corrected node loss is
\begin{equation}
\label{eq:v8-corrected-node-loss}
  \ell_v^{\mathrm{corr}}
  =
  \widehat\ell_v
  +
  \frac{R_v}{\pi_v}
  \left(\ell_v^\ast-\widehat\ell_v\right).
\end{equation}
If \(R_v=0\), the node contributes the proxy loss. If the node is sampled, the
observed oracle--proxy difference is added back with its inverse-propensity
weight. With correct logged propensities,
\(\mathbb E[\ell_v^{\mathrm{corr}}\mid T]=\ell_v^\ast\).

Aggregating the corrected node losses gives the practical local-law channel:
\[
  L_{\mathrm{law}}^{\mathrm{corr}}(g,T)
  =
  \sum_{v\in V(T)}\gamma^{d(v)}\ell_v^{\mathrm{corr}} .
\]
It estimates \(L_{\mathrm{law}}^\ast(g,T)\); it is not a replacement population
target. In finite data, the objective actually optimized or reported is the
plug-in version
\begin{equation}
\label{eq:v8-corrected-objective}
  J_\lambda^{\mathrm{corr}}(g;T)
  =
  (1-\lambda)L_{\mathrm{root}}(g)
  +
  \lambda L_{\mathrm{law}}^{\mathrm{corr}}(g,T).
\end{equation}

The remaining bookkeeping is design-based estimation. Coverage changes the
variance of the root and local-law estimates, not the population objective.
Explicit channel weights are an equivalent interface to \(\lambda\) mode, but
the two interfaces should not be supplied at the same time. Proxy-oracle gaps,
propensity residuals, sampling standard errors, and clipping envelopes are
diagnostics for the corrected channel, not additive loss terms. The full
Horvitz--Thompson algebra, explicit-weight form, and diagnostic decomposition
are in Appendix~\ref{app:v8-objective-ipw}.

\subsection{Putting the Pieces Together}

The full C-TreePO picture is therefore short enough to state without the
estimation details. First, choose a realized C-Tree for a document, so every
leaf, internal node, and root has a represented span. Second, learn one state
map \(g\) that is used throughout that tree. Third, train \(g\) against two
channels: a root channel for the document-level target and a local-law channel
for C1/C2/C3 validity at the nodes. Fourth, because \(f^\ast\) is usually
sparse, use \(\widehat f\) for dense local training and sampled \(f^\ast\)
node calls to audit and correct the local-law channel.

Section~\ref{sec:v8-ctree-math} supplies the semantic reason this stack is
meaningful: if the local laws hold exactly, node and root states can replace
their raw spans for oracle-compatible readouts. This section supplies the
optimization and observation layer: \(J_\lambda^\ast\) is the oracle objective,
and \(J_\lambda^{\mathrm{corr}}\) is its proxy-corrected finite-data version.
The correction estimates the oracle local-law channel; it is not a new target.

The Manifesto study below should be read as an LLM instantiation of this stack.
Expert-survey scores anchor the root channel, teacher traces and proxy scores
shape the learned state map \(g\), and the stronger node-level \(f^\ast\) audit
is the next measurement layer.
